% Part: first-order-logic
% Chapter: completeness
% Section: outline

\documentclass[../../../include/open-logic-section]{subfiles}

\begin{document}

\iftag{FOL}
      {\olfileid{fol}{com}{out}}
      {\olfileid{pl}{com}{out}}

\olsection{నిరూపణ రూపురేఖ}

సంపూర్ణతా సిద్ధాంతపు నిరూపణ కొంత సంక్లిష్టమైనది; దాన్ని మొదటిసారి
చదివేటప్పుడు దారి తప్పడం సులభం. అందుకే నిరూపణకు ఒక రూపురేఖ
వేసుకుందాం. మొదటి అడుగు దృక్కోణాన్ని మార్చడం; అది నిరూపణకు ఒక మార్గాన్ని
చూడనిస్తుంది. సంపూర్ణతను ``$\Gamma \Entails !A$ అయినప్పుడల్లా
$\Gamma \Proves !A$'' అని ఆలోచించినప్పుడు ఒక ఆలోచన రావడమే కష్టంగా
ఉండవచ్చు: $\Gamma \Proves !A$ అని చూపాలంటే మనం !!a{derivation}ను
కనుగొనాలి; కానీ $\Gamma \Entails !A$ అనే పరికల్పన అందుకు ఏ విధంగానూ
సహాయపడుతున్నట్లు కనిపించదు. కొన్ని నిరూపణ వ్యవస్థలకు నేరుగా
!!a{derivation}ను నిర్మించవచ్చు; అయితే మనం కొద్దిగా భిన్నమైన పద్ధతిని
తీసుకుంటాం. కావలసిన దృక్కోణ మార్పు ఇది: సంపూర్ణతను ``$\Gamma$
అవైరుధ్యమైతే, అది సంతృప్తిపరచదగినది'' అని కూడా సూత్రీకరించవచ్చు.
~$\Gamma$లోని సమాచారాన్ని, అది అవైరుధ్యమనే పరికల్పనతో కలిపి,
$\Gamma$లోని ప్రతి \iftag{FOL}{!!{sentence}}{!!{formula}}ను
సంతృప్తిపరచే \iftag{FOL}{!!a{structure}}{!!a{valuation}}ను మనం
నిర్మించగలమేమో. ఎంతైనా, మనం వెతుకుతున్న
\iftag{FOL}{!!{structure}}{!!{valuation}} ఎలాంటిదో తెలుసు:
$\Gamma$ దాన్ని ఎలా వర్ణిస్తుందో సరిగ్గా అలాంటిదే!{}

$\Gamma$లో \iftag{FOL}{పరమాణు !!{sentence}s}{!!{propositional variable}s}
మాత్రమే ఉంటే, దానికి ఒక నమూనాను నిర్మించడం సులభం.
\iftag{FOL}{ఆ పరమాణు !!{sentence}s అన్నీ $\Atom{P}{a_1,\dots,a_n}$
  రూపంలో ఉన్నాయని, ఇక్కడ $a_i$లు !!{constant}s అని అనుకుందాం.}{}
మనకు చేయవలసిందల్లా
\iftag{FOL}{%
  !!a{domain}~$\Domain{M}$ను, $\Sat{M}{\Atom{P}{a_1,\ldots,a_n}}$
  అయ్యేలా~$P$కు ఒక నిర్దేశాన్ని ఏర్పరచడమే. అది పెద్ద కష్టం కాదు:
  $\Domain{M} = \Nat$గా, $\Assign{\Obj c_i}{M} = i$గా ఉంచి, ప్రతి
  $\Atom{P}{a_1,\ldots,a_n} \in \Gamma$కు $\tuple{k_1,
    \dots, k_n}$ను $\Assign{P}{M}$లో ఉంచండి; ఇక్కడ $k_i$ అనేది
  స్థిరాంక సంకేతం~$a_i$ యొక్క సూచిక (అంటే, $a_i \ident \Obj c_{k_i}$).}{%
  అన్ని $p \in \Gamma$లకు $\pSat{v}{p}$ అయ్యే
  !!a{valuation}~$\pAssign{v}$ను ఏర్పరచడమే. $p \in \Gamma$ అయినప్పుడు,
  అప్పుడు మాత్రమే $\pAssign{v}(p) = \True$గా ఉంచుదాం.}

ఇప్పుడు~$\Gamma$లో $\lnot !B$ అనే ఏదో !!{formula} ఉందనీ, $!B$
పరమాణు సూత్రమనీ అనుకుందాం. $\iftag{FOL}{\Struct{M}}{\pAssign{v}}$
నిర్మాణం $\lnot !B$ను సత్యం చేయగల అవకాశానికి ఆటంకం కలిగిస్తుందేమోనని
ఆందోళన పడవచ్చు. ఇక్కడే~$\Gamma$ యొక్క అవైరుధ్యం ఉపయోగపడుతుంది:
$\lnot !B \in \Gamma$ అయితే $!B \notin \Gamma$; లేకపోతే $\Gamma$
వైరుధ్యభరితమవుతుంది. $!B \notin \Gamma$ అయితే, మన
$\iftag{FOL}{\Struct{M}}{\pAssign{v}}$ నిర్మాణం ప్రకారం
$\iftag{FOL}{\Sat/{M}}{\pSat/{v}}{!B}$; అందువల్ల
$\iftag{FOL}{\Sat{M}}{\pSat{v}}{\lnot !B}$. ఇప్పటివరకు అంతా సవ్యంగానే
ఉంది.

$\Gamma$లో సంక్లిష్టమైన, పరమాణుకాని సూత్రాలు ఉంటే ఏమవుతుంది? దానిలో
$!A \land !B$ ఉందనుకుందాం. దాన్ని సత్యం చేయడానికి $!A$, $!B$ రెండూ
~$\Gamma$లో ఉన్నట్లుగా ముందుకు సాగాలి. $!A \lor !B \in \Gamma$ అయితే,
వాటిలో కనీసం ఒకదాన్ని సత్యం చేయాలి; అంటే, వాటిలో ఒకటి~$\Gamma$లో
ఉన్నట్లుగా ముందుకు సాగాలి.

ఇది కింది ఆలోచనను సూచిస్తుంది: (a)~ఫలిత సమితిని అవైరుధ్యంగా ఉంచేలా,
(b)~సంభవించే ప్రతి !!{sentence}~$!A$కు ఫలిత సమితిలో $!A$ లేదా
$\lnot !A$ ఉండేలా, (c)~సమితిలో $!A \land !B$ ఉన్నప్పుడల్లా $!A$,
$!B$ రెండూ కూడా ఉండేలా, $!A \lor !B$ ఉంటే $!A$ లేదా $!B$లో కనీసం
ఒకటి కూడా ఉండేలా, తదితరంగా, మరిన్ని !!{formula}sను~$\Gamma$కు
చేర్చుతాం. \sourcecorrection{OLTECOM-001}{ఆంగ్ల మూలం (b)లో ప్రతి
``పరమాణు'' వాక్యానికి ఎంపిక ఉండాలని చెబుతుంది; కానీ వెంటనే ఇచ్చిన
సంపూర్ణ సమితి నిర్వచనం, తరువాతి నిర్మాణం ప్రతి వాక్యానికి ఆ ఎంపికను
కోరుతాయి. అందుకే ఇక్కడ పరమాణు అనే అనవసర పరిమితిని తొలగించాం.}
దీన్ని (అవసరమైతే అంతులేని కాలం) కొనసాగిస్తాం. ఇలా చేర్చిన అన్ని
!!{formula}s సమితిని~$\Gamma^*$ అందాం. అప్పుడు పై నిర్మాణం మనకు
\iftag{FOL}{!!a{structure}~$\Struct{M}$}{!!a{valuation}~$\pAssign{v}$}ను
ఇస్తుంది; అది~$\Gamma^*$లోని అన్ని వాక్యాలను సంతృప్తిపరుస్తుందని
ఆగమనంతో నిరూపించవచ్చు. $\Gamma \subseteq \Gamma^*$ కనుక, అది
~$\Gamma$లోని అన్ని వాక్యాలను కూడా సంతృప్తిపరుస్తుంది. (a), (b)లకు
హామీ ఇవ్వడమే చాలుతుందని తేలుతుంది. (b) నెరవేర్చే వాక్యాల సమితిని
\emph{సంపూర్ణమైనది} అంటారు. కాబట్టి అవైరుధ్య సమితి~$\Gamma$ను
అవైరుధ్యమైన, సంపూర్ణమైన సమితి~$\Gamma^*$గా విస్తరించడం మన పని.

\iftag{FOL}{%
ఈ ప్రణాళికలో ఒక చిక్కు ఉంది: $\lexists[x][!A(x)] \in \Gamma$ అయితే,
ఏదో ఒక !!{constant}~$c$ను ఎంచుకొని ఈ ప్రక్రియలో $!A(c)$ను చేర్చాలని
ఆశిస్తాం. కానీ అది ఎల్లప్పుడూ చేయగలమని ఎలా తెలుసు? మన భాషలో కొన్ని
!!{constant}s మాత్రమే ఉండవచ్చు; వాటిలో ఒక్కొక్కదానికీ $\lnot !A(c)
\in \Gamma$ ఉండవచ్చు. అప్పుడు $!A(c)$ను కూడా చేర్చలేం, ఎందుకంటే అది
సమితిని వైరుధ్యభరితం చేస్తుంది; పైగా $\Struct{M}$ అనేది $!A(c)$ను
సత్యం చేయాలా, $\lnot !A(c)$ను సత్యం చేయాలా అనేది మనకు తెలియదు.
అంతేకాక, $\Gamma$లో స్థిరాంక సంకేతాలే లేని భాషలోని వాక్యాలు మాత్రమే
ఉండవచ్చు (ఉదాహరణకు, సమితి సిద్ధాంతపు భాష).

ఈ సమస్యకు పరిష్కారం ఆరంభంలోనే అనంత సంఖ్యలో స్థిరాంకాలను, అలాగే వాటిని
పరిమాణీకరణికాలతో సరైన విధంగా అనుసంధానించే వాక్యాలను చేర్చడమే. (ఇది
వైరుధ్యాన్ని ప్రవేశపెట్టదని తప్పకుండా నిర్ధారించాలి.)

పరమాణు వాక్యాల్లో !!{constant}s మాత్రమే ఉన్నప్పుడు మన మొదటి నిర్మాణం
చక్కగా పనిచేస్తుంది. కానీ భాషలో !!{function}s కూడా ఉండవచ్చు. అప్పుడు
అన్నీ పనిచేసేలా ఆ !!{function}sకు నిర్దేశించడానికి~$\Nat$ మీద సరైన
ప్రమేయాలను కనుగొనడం కష్టమవచ్చు. అందుకే మరొక ఉపాయం: $\Obj c_i$ను
అర్థం చేసుకోవడానికి $i$ని ఉపయోగించకుండా, !!{constant}s సమితినే
ప్రదేశంగా తీసుకుందాం. అప్పుడు $\Struct M$ ప్రతి !!{constant}ను దానికే
నిర్దేశించగలదు: $\Assign{\Obj c_i}{M} = \Obj c_i$. మరి ఇంకా ముందుకు
ఎందుకు వెళ్లకూడదు? భాషలోని అన్ని
\emph{పదాలనే}~$\Domain{M}$గా తీసుకుందాం!{} అలా చేస్తే,
!!{function}sకు (పదాలను ఆర్గ్యుమెంట్లుగా తీసుకొని పదాలను విలువలుగా
ఇచ్చే) ప్రమేయాలను నిర్దేశించే స్పష్టమైన పద్ధతి ఉంది: $n$ పదాలు
$t_1$, \dots,~$t_n$ను ఇన్‌పుట్‌గా తీసుకొని $\Obj f^n_i(t_1, \dots,
t_n)$ అనే పదాన్ని విలువగా ఇచ్చే ప్రమేయాన్ని !!{function}~$\Obj f^n_i$కు
నిర్దేశిస్తాం.

చివరి చిక్కు~$\eq$ను ఏం చేయాలనేది. !!{predicate}~$\eq$కు స్థిరమైన
అర్థవివరణ ఉంది: $\Value{t}{M} = \Value{t'}{M}$ అయినప్పుడు, అప్పుడు
మాత్రమే $\Sat{M}{\eq[t][t']}$. ఒక పదం~$t$ యొక్క !!{value}ను $t$
అదే అయ్యేలా ఏర్పాటు చేస్తే, $t$, $t'$ ఒకే పదం అయినప్పుడు తప్ప ఈ
!!{structure} $\eq[t][t']$ రూపంలోని \emph{ఏ} వాక్యాన్నీ సత్యం చేయదు. ఇది
సమస్యే; ఎందుకంటే !!{function}s గల భాషలోని దాదాపు ప్రతి ఆసక్తికరమైన
సిద్ధాంతంలో, $t$, $t'$ ఒకే పదం కాకపోయినా $\eq[t][t']$ అనే వాక్యాలు
సిద్ధాంతాలుగా ఉంటాయి (ఉదాహరణకు, అంకగణిత సిద్ధాంతాల్లో:
$\eq[(\Obj 0+ \Obj 0)][\Obj 0]$). ఈ సమస్యను పరిష్కరించడానికి
~$\Struct M$ యొక్క ప్రదేశాన్ని మారుస్తాం: $\Domain{M}$లో పదాలను
వస్తువులుగా ఉపయోగించకుండా పదాల సమితులను ఉపయోగిస్తాం; ఒక్కొక్క సమితిలో
~$\Gamma$లోని వాక్యాలు సమానంగా ఉండాలని కోరే పదాలన్నీ ఉంటాయి.
ఉదాహరణకు, $\Gamma$ అంకగణిత సిద్ధాంతమైతే, ఈ సమితుల్లో ఒకదానిలో
$\Obj 0$, $(\Obj 0 + \Obj 0)$, $(\Obj 0 \times \Obj 0)$ మొదలైనవి
ఉంటాయి. ఈ సమితినే $\Obj 0$కు నిర్దేశిస్తాం; దానిలోని అన్ని పదాలకూ,
ఉదాహరణకు $(\Obj 0 + \Obj 0)$కు కూడా, ఇదే సమితి విలువ అవుతుందని
తేలుతుంది. అందువల్ల సవరించిన ఈ !!{structure}లో
$\eq[(\Obj 0+ \Obj 0)][\Obj 0]$ సత్యమవుతుంది.}{}

ఇప్పుడు మనం చేయబోయేది ఇదే. మొదట !!{complete} అవైరుధ్య సమితుల ధర్మాలను
పరిశీలిస్తాం. ప్రత్యేకంగా, !!a{complete} అవైరుధ్య సమితిలో $!A \land
!B$ ఉండటం, $!A$,~$!B$ రెండూ ఉండటానికి సరిసమానమని; $!A \lor !B$
ఉండటం వాటిలో కనీసం ఒకటి ఉండటానికి సరిసమానమని; తదితర విషయాలను
నిరూపిస్తాం (\olref[ccs]{prop:ccs}).\iftag{FOL}{ తరువాత ``సంతృప్తీకృత''
  వాక్యాల సమితులను నిర్వచించి పరిశీలిస్తాం. ప్రతి పరిమాణీకృత
  !!{sentence}ను దాని రూపాలతో అనుసంధానించే షరతులను కలిగిన సమితినే
  సంతృప్తీకృత సమితి అంటాం (\olref[hen]{defn:henkin-exp}). ఏ
  అవైరుధ్య సమితి~$\Gamma$నైనా ఎల్లప్పుడూ ఒక సంతృప్తీకృత
  సమితి~$\Gamma'$గా విస్తరించవచ్చని చూపుతాం
  (\olref[hen]{lem:henkin}). ఒక సమితి అవైరుధ్యమైనది, సంతృప్తీకృతమైనది,
  !!{complete} అయినప్పుడు, దానిలో
  \iftag{prvEx}{$\lexists[x][!A(x)]$ ఉండటం, ఏదో ఒక సంవృత పదం~$t$కు
    $!A(t)$ ఉండటానికి సరిసమానం}{}\iftag{defEx,defAll}{}{
    మరియు }\iftag{prvAll}{$\lforall[x][!A(x)]$ ఉండటం, అన్ని సంవృత
    పదాలు~$t$కు $!A(t)$ ఉండటానికి సరిసమానం}{} అనే ధర్మం కూడా ఉంటుంది
  (\olref[hen]{prop:saturated-instances}).}{} తరువాత
\iftag{FOL}{సంతృప్తీకృత}{} అవైరుధ్య సమితి
~$\iftag{FOL}{\Gamma'}{\Gamma}$ను తీసుకొని, దాన్ని
\iftag{FOL}{సంతృప్తీకృతమైన, అవైరుధ్యమైన,}{అవైరుధ్యమైన మరియు}
!!{complete} సమితి~$\Gamma^*$గా విస్తరించవచ్చని చూపుతాం
(\olref[lin]{lem:lindenbaum}). ఈ $\Gamma^*$ను మన
\iftag{FOL}{పద నమూనా~$\Struct M(\Gamma^*)$}{!!{valuation}~$\pAssign
v(\Gamma^*)$}ను నిర్వచించడానికి ఉపయోగిస్తాం. \iftag{FOL}{పద నమూనా
  ప్రదేశం సంవృత పదాల సమితి; దాని !!{predicate}s అర్థవివరణను}{ఆ
  మూల్యనిర్ణయాన్ని}~$\Gamma^*$లోని \iftag{FOL}{పరమాణు
  !!{sentence}s}{!!{propositional
    variable}s} నిర్ణయిస్తాయి
(\olref[mod]{defn:termmodel}). \iftag{FOL}{సంతృప్తీకృత,}{} సంపూర్ణ
అవైరుధ్య సమితుల ధర్మాలను ఉపయోగించి నిజంగానే
$\iftag{FOL}{\Sat{M(\Gamma^*)}}{\pSat{v(\Gamma^*)}}{!A}$ అయినప్పుడు,
అప్పుడు మాత్రమే $!A \in \Gamma^*$ అని
(\olref[mod]{lem:truth}), ముఖ్యంగా
$\iftag{FOL}{\Sat{M(\Gamma^*)}}{\pSat{v(\Gamma^*)}}{\Gamma}$ అని
చూపుతాం. \iftag{FOL}{చివరగా, $\Gamma$లో~$\eq$ కూడా ఉంటే పద నమూనాను
  ఎలా నిర్వచించాలో పరిశీలించి
  (\olref[ide]{defn:term-model-factor}), అది~$\Gamma^*$ను
  సంతృప్తిపరుస్తుందని చూపుతాం (\olref[ide]{lem:truth}).}{}

\end{document}
